%! Author = Saeid Yazdani
%! Date = 9/6/2024

\section{Software Design Approach}
\label{sec:software-design-approach}

In the context of embedded systems and measurement equipment, two
approaches are often considered: the \emph{Operating System} approach and the
\emph{Bare-Metal} approach.

As of recent, the industry has increasingly shifted toward \acs{O/S}-based
system software, particularly in more complex and feature-rich equipment
such as high-end oscilloscopes and signal analyzers.

The use of \acs{SoC}s allows for the integration of powerful \acs{FPGA}s and
\acs{CPU}s 
capable of running operating systems like \emph{Linux}, which provide a
rich user interface, multitasking, and connectivity
options such as Ethernet, Wi-Fi, and USB.
This trend is driven by the demand for advanced features like touchscreen
interfaces, remote access, data logging, and the ability to run sophisticated
built-in analysis software.

However, the premium look and feel of the user interface and software
capabilities of such systems are associated with a set of challenges, of which:

\begin{itemize}[topsep=8pt,itemsep=4pt,partopsep=4pt, parsep=4pt]
    \item \textbf{Development Complexity:}
    The requirement for managing multiple layers such as the
    \acs{O/S} kernel, drivers, libraries, user-space applications and
    all the interdependencies between the layers.

    \item \textbf{Software System Reliability:}
    \acs{O/S}-Based systems, particularly those handling both user interfaces
    and hardware control may face reliability issues such as kernel panics
    and software crashes which are unacceptable and often
    considered a safety issue for equipment such as power supplies.

    \item \textbf{Security Concerns:}
    Vulnerabilities in the \acs{O/S}, libraries, or custom software can lead to
    unauthorized access or data breaches if the equipment is networked or offers
    remote access.

    \item \textbf{Startup and Task Latencies:}
    The equipment may require fast startup times and immediate availability.
    Booting a full \acs{O/S}, especially with \acs{GUI} and other services,
    introduces delay.
    These operating systems are often not optimized for realtime tasks
    due to their scheduling strategy.

\end{itemize}

On the contrary, a bare-metal system allows for direct control over hardware
with no operating system overhead that simplifies development
as there is no need to manage complex \acs{O/S}-related tasks such as
multitasking or memory management.
Minimal latency and predictable behavior can be easily achieved by direct
interaction with the hardware.

However, bare-metal development can become effort-intensive for complex systems
requiring advanced features (e.g., networking, GUIs) that would be easier to
implement on an \acs{O/S}.
Therefore, the choice between \acs{O/S}-based or bare-metal approach is tightly
coupled with the expected functionality of the equipment, the complexity of
the data handling and processing as well as the size and skill of
the development team.
For these reasons, bare-metal software development is the selected approach
for the proposed \acs{DPS}.